\section{Component constants and the geometric target representation}

% BEGIN A77G1
\paragraph{A77G1. Correcting all components by a unit scalar.}
Consider the right action $x\mapsto xg$ of a finite group $G$ on
components. An adelic point $h$ has component $x(h)$, and each
lift $k$ of $g$ takes it to $x(h)g$. Lifts exist and multiply.
Let $r(f,h)$ be a reading, $\rho$ a left action on the forms,
$\mu$ a unitary character on the adelic points, and $s$ a character
of $G$. Assume the two geometric move laws. For every permitted lift,
\[
 r(\rho(g)f,h)=z(x(h),g)r(f,hk),\quad |z(x,g)|=1 ,
\]
and when $x(h)g=x(h)$ the equality also holds with multipliers
\[
 \mu(h) r(\rho(g)f,h)=s(g)\mu(hk)r(f,hk).
\]
Then there are unit constants $c(x)$ making the latter equality
hold for every move after multiplying each side by $c$ at its
respective component.

Call $x$ active if some $r(f,h)$ there is nonzero. The first
move law applied to $\rho(g^{-1})f$ shows that activity is
invariant in both directions. At an active $x$, applying both
laws to a lift of the identity and cancelling a nonzero reading
gives $\mu(k)=z(x,1)$. Multiplying two different lifts of $g$
on the right by the same lift of $g^{-1}$ now shows that their
$\mu$-values agree. Set
\[
 w(x,g)=z(x,g)/(s(g)\mu(k))
\]
on active components and set $w=1$ elsewhere. Since $G$ is
finite, $s(g)$ is a root of unity. Every $w$ has absolute value
one. The move formula gives
$\mu(h)r(\rho(g)f,h)=w(x,g)s(g)\mu(hk)r(f,hk)$ everywhere,
including the identically zero components. On an active component
choose a nonzero reading of $f$ and apply this equality to
$\rho((gg')^{-1})f$ once with the product lift and once with its
two factors. Cancellation gives
$w(x,gg')=w(x,g)w(xg,g')$. Applying the stabilizer formula in
the same way gives $w(x,g)=1$ if $xg=x$. Both identities hold
also on the other orbits. On each orbit fix a representative $x_*$,
and define $c(x_*g)=w(x_*,g)$. The stabilizer identity makes this
well-defined and the product identity gives
$c(xg)=c(x)w(x,g)$. This is the required normalization even on
components where every reading vanishes.
% END A77G1

% BEGIN A77G2
\paragraph{A77G2. Installing the base solution without changing its phase.}
Use the tied core, $R$ and the reference rows of P77R2. At the
distinguished embedding take the solution
$O_i.\mathrm{OneMapOf}(c_i,1)$ supplied by P77R3.
The level law, the ordered equality $T.\mathrm{pb}=R.C.P.\mathrm{pb}$,
$\mathrm{TieIn}$ and the reference pins suffice for the other good
embeddings as well. One can prescribe a single
unit family $c_E$ on the core which has $c_i$ at the distinguished
embedding and satisfies the native predicate $\mathrm{Fired}(1)$.

For the character $s$, each factor of the product of residue
determinant characters at $g$ has its value in the small field.
Indeed lift its residue determinant to an integral unit at that
place, insert it there and insert one at all other places. The
reference character tie recovers the residue value, up to inversion.
The product therefore has a unique preimage $y_g$ in the small
field. Injectivity at the distinguished embedding gives
$y_1=1$ and $y_{gg'}=y_g y_{g'}$. Put $s(g)=\sigma(y_g)$ at
the embedding $\sigma$ in question. In the first move law,
$r$ is the reading of a cusp section in its dictionary frame;
$z(x,g)$ is the unit from comparing the moved period frames.
The dictionary and adelic pins in $\mathrm{TieIn}$ identify the
lattice frame, the finite level, and the rational mover at that
component. The reference row's lift pin supplies a lift for each $g$,
and its action pin identifies the induced change of the marking.
These give the right component action and its raw reading formula:
right multiplication by the lift corresponds to pullback on the
marked section, up to the unit comparing the normalized period
generators. The framed readout comparison in the core gives this
same unit for every section and every lift of that $g$ at the
component. On a stabilized component the self-component comparison
can be made in the original chart. The determinant of the lift has
the value prescribed by the tied row character, giving precisely
the second move formula with $s(g)\mu(hk)$ in A77G1. These uses of
the core and reference pins all precede the target-coordinate action. A77G1 gives
the unit constants and the required $\mathrm{EqAt}$ identity
for every lift, with this same $y_g$.

At the distinguished embedding use $c_i$ itself. Its equivariance
is the first field of the supplied base solution and its modulus
is one by the row law. Use the preceding construction at each
other good embedding where the row space is the specified target.
Set the remaining entries, the zeroth row and unused rows equal
to one. This defines $c_E$ before evaluating any section. Substituting
it in the core preserves all fields of $\mathrm{TieIn}$: the only
new core law is the unit modulus of these entries. Reference pins
have no such constants among their entries. By construction the
base solution has exactly its old constants and phase one; at
every required good embedding we have the second, equivariance,
clause of $\mathrm{Fired}$.
% END A77G2

% BEGIN A77G3
\paragraph{A77G3. Constructing the geometric half on the corrected coordinates.}
Start with $\mathrm{PreChar}$, $\mathrm{TieIn}$, the reference
rows and the solution $\mathrm{Fired}(1)$ of A77G2, on the same
core and real record. For row $i$ put $C_i=\mathrm{Cu}(T,\kappa,i)$,
the original cusp space over the core field $K$. It is finite
dimensional by the core laws. Choose the basis in the source's
coordinate map, of size $n_i=\dim_K C_i$. In $L^{n_i}$ the
$K$-form is the image of $C_i$ in these coordinates. The seminorm
at a geometric frame is the seminorm of its inverse coordinate
section there. At $\sigma:L\to\mathbb C$ define the realization
$J_\sigma$ by extending the images of the basis vectors under
the original $j_{\rm Top,\sigma}$ semilinearly. Thus for every $f$,
\[
     J_\sigma(\operatorname{coord} f)=j_{\rm Top,\sigma}(f).
\]
This definition includes the conditional evaluation branches
of $j_{\rm Top}$. When the evaluation lift exists, its values
on the marked vector are the section readings by definition;
when it fails the corresponding zero clause applies. The same
basis extension gives the analogous conjugate evaluation clauses.
The coordinate image, inverse norm test and these evaluation
identities give the weight and frame tie and its conjugation law.

The action stored on these coordinates is prescribed too.
If the marked level action preserves $C_i$, restrict that linear
action to $C_i$. The group unit and product laws follow by
evaluating the unit and the composite level changes on each
section, using its naturality. In the chosen basis extend its
matrices from $K$ to $L$. Their products are the matrices of the
composites; the resulting action has exactly the prescribed effect
on each coordinate cusp section. Conversely, a representation
with this effect forces preservation since its image lies in
$C_i$. In the alternative branch the definition of
$\mathrm{geoRho}$ therefore gives the identity action, which is
also the required fallback law. The determinant character in the
coordinate record is the one defined by $g\mapsto y_g$ in A77G2
and extension to $L$.

Replace the two target fields of $R$ by these $n_i$ and these
forms with their action, producing $R_{\rm geo}$.
The reference tuple data use the original model, marked level,
theta functions and definite inputs. Each is identical after this
replacement. Thus the reference pins and both reading identities
in $\mathrm{Fired}(1)$ hold for the copied rows. The same check on
the model, family, chain and theta fields gives $\mathrm{PreChar}$
for $R_{\rm geo}$ and the determinant calculation gives its
group character law. The tie, conjugation, reference pins, action
and fired identities supply all fields of
$\mathrm{GeoHalf}(T,R_{\rm geo},1)$: equivariance at the other
embeddings is the second fired identity, while the comparison
and the guarded phase clause use the first. This constructs a
new set of coordinates on the fixed analytic model and supplies
the geometry on precisely those coordinates.
% END A77G3
